\chapter*{怎样读这一卷}
\addcontentsline{toc}{chapter}{怎样读这一卷}

本卷写前二〇六年至二二〇年的西汉、新、东汉与汉魏禅代。秦末只用来说明开国条件，二二〇年则是本卷必须读完的制度断点；不要把它当作三国故事已经开始的注脚。诸侯国、郡县、边郡与朝廷都在同一时间线上行动，地方社会和边地不是长安、雒阳叙事的背景。

先随年份、年号和政权转换阅读事件簇，再回看专题怎样改变这条年序的压力：征发、财政、边疆、继承和地方网络如何相互牵动。相邻事件不等于互为原因；请分别辨认较长期的条件、眼前的触发、使行动成为可能的资源与道路，以及直接结果如何留下下一阶段的难题。

《史记》《汉书》《后汉书》《后汉纪》与《资治通鉴》提供编年的骨架；《食货志》、盐铁材料和居延、肩水金关、悬泉置简牍让我们追问制度在何处可见。它们都不是全社会的透明记录：诏令不等于已经实行，边塞文书不等于全国，军数、户口、路线和人物动机也常只能写到材料允许的程度。

页中的信息框是阅读的停靠点：\texttt{event} 标出首次深描的事件，\texttt{turningpoint} 提醒后续路径何处转向，\texttt{causalchain} 分开条件、触发与后果；\texttt{textwindow} 只解释一小段原文的语境；\texttt{sourcenote} 和 \texttt{debate} 则说明证据能证实什么、争议尚留在哪里。把这些框同前后的叙事一起读，而不要把它们当作结论的替代品。
